\documentclass[10pt]{article}
\usepackage[a4paper,margin=26mm]{geometry}
\usepackage[T1]{fontenc}
\usepackage{lmodern}
\usepackage{amsmath,amssymb}
\emergencystretch=2em
\title{Complement is an isometric involution\\Proof of Conjecture 00000006406}
\author{ziangni-sys}
\date{4 October 2026}
\begin{document}
\maketitle
\section*{Set identity and distance preservation}
Fix a universe $X$ and write $A^c=X\setminus A$ for complement.
All complements are relative to this same universe.
For arbitrary subsets $A,B\subseteq X$, their symmetric difference is
$A\triangle B=(A\setminus B)\cup(B\setminus A)$.
De Morgan's identities are
\[
(A\cup B)^c=A^c\cap B^c,\qquad
(A\cap B)^c=A^c\cup B^c.
\]
Together with double complementation they give
$A^c\setminus B^c=A^c\cap B=B\setminus A$. Therefore
\[
A^c\triangle B^c=(B\setminus A)\cup(A\setminus B)
=A\triangle B.
\]
This proves preservation of every symmetric-difference distance:
if $d(A,B)=W(A\triangle B)$ for any fixed set-size functional $W$, then
\[
d(A^c,B^c)=W(A^c\triangle B^c)
=W(A\triangle B)=d(A,B).
\]
Whenever this distance is a metric, complement is an isometry.
The equality itself requires no finiteness or positivity assumption.
Examples include counting distance on a finite universe, and
measure distance on a complement-closed measurable family
(with sets identified modulo null sets when necessary).

\section*{Involution and its composition}
Let $C(A)=A^c$. Pointwise double complementation gives
$C(C(A))=A$. Thus $C$ is an involution and is bijective, with inverse $C$.
Moreover $C\circ C$ is the identity map and is itself an involution.
This proves both the usual involution assertion and the literal wording
about the composition of complements.
For a restricted family, complement is considered on a family closed
under this operation; on the full power set this condition is automatic.

\section*{Formalization and a concrete metric}
The Lean project proves the two De Morgan identities pointwise, the
set-difference identity and the symmetric-difference equality above.
It proves distance preservation for an arbitrary functional $W$, then
an actual Mathlib \texttt{Isometry} theorem for every pseudometric whose
distance has the stated symmetric-difference form.
The hypothesis specifies what the difference distance is; it does not
assume that complement preserves it.
The final theorem \texttt{conjecture\_6406} collects isometry, both
involution statements, and both De Morgan identities.

Additionally, finite subsets are represented by their Boolean
characteristic words. Mathlib's actual Hamming metric counts entries at
which two such words differ, exactly the cardinality of their symmetric
difference. Bitwise Boolean negation represents complement.
The unconditional theorem \texttt{concrete\_hamming\_isometry} proves
it is an isometry for every finite universe, and double negation is
proved separately. All proofs are exact; no auxiliary computation is used.
\end{document}
